```latex
\documentclass[11pt]{article}

\usepackage{amsmath,amssymb,amsthm}
\usepackage[colorlinks=true,linkcolor=blue,citecolor=blue,urlcolor=blue]{hyperref}

\newtheorem{theorem}{Theorem}[section]
\newtheorem{proposition}[theorem]{Proposition}
\newtheorem{lemma}[theorem]{Lemma}
\newtheorem{corollary}[theorem]{Corollary}
\theoremstyle{definition}
\newtheorem{definition}[theorem]{Definition}
\theoremstyle{remark}
\newtheorem{remark}[theorem]{Remark}

\newcommand{\cC}{\mathbb{C}}
\newcommand{\cN}{\mathcal N}
\newcommand{\cU}{\mathcal U}
\newcommand{\cA}{\mathcal A}
\newcommand{\cJ}{\mathcal J}
\newcommand{\Symp}{\operatorname{Symp}}
\newcommand{\Aux}{\operatorname{Aux}}
\newcommand{\AK}{\operatorname{AK}}
\newcommand{\id}{\mathrm{id}}

\begin{document}

\title{Holomorphic branch coordinates near symplectic normal crossings divisors}
\author{}
\date{}
\maketitle

\begin{abstract}
We prove that every symplectic simple normal crossings divisor admits, after the
usual cohomologous regularizing deformation, a tame almost complex structure
with locally holomorphic defining equations for the branches.
\end{abstract}

\section{Statement}

Let \((M,\omega)\) be a closed symplectic \(2n\)-manifold and let
\[
D=\bigcup_{a=1}^{N}D_a
\]
be a symplectic simple normal crossings divisor. For
\(I\subseteq\{1,\ldots,N\}\), write
\[
D_I=\bigcap_{a\in I}D_a .
\]

For a complex-valued smooth function \(f\) on an almost complex manifold
\((M,J)\), we say that \(f\) is \(J\)-holomorphic if
\[
df\circ J=i\,df .
\]

\begin{theorem}\label{thm:main}
There is a smooth path \((\omega_t)_{0\leq t\leq1}\) of symplectic forms on \(M\)
such that
\begin{enumerate}
\item \([\omega_t]=[\omega]\in H^2(M;\mathbb R)\) for every \(t\);
\item \(D\) is a symplectic simple normal crossings divisor for every
\(\omega_t\);
\item there is an \(\omega_1\)-tame almost complex structure \(J\) with the
following property: for every \(p\in M\), there are a neighborhood \(U\ni p\),
smooth complex coordinates
\[
(z_1,\ldots,z_n):U\longrightarrow \cC^n,
\]
and an integer \(0\leq r\leq n\) such that
\[
dz_a\circ J=i\,dz_a,\qquad 1\leq a\leq r,
\]
and
\[
D\cap U=\{z_1\cdots z_r=0\}.
\]
\end{enumerate}
Here \(r=0\) is used at points of \(M\setminus D\); the empty product is \(1\).
\end{theorem}

Since \(D\) is a proper divisor, at a point of depth \(r\) one necessarily has
\(r<n\) unless one is considering the empty divisor.

\section{Regularization}

We use the standard terminology of symplectic normal crossings divisors,
\(\omega\)-regularizations, and almost K\"ahler data. Let
\(\Symp^+(M,D)\) denote the space of symplectic forms for which \(D\) is a
positively symplectic SC divisor, and let \(\Aux(M,D)\) denote the space of
pairs \((\omega,\mathcal R)\), where \(\mathcal R\) is an
\(\omega\)-regularization of \(D\).

\begin{theorem}[Tehrani--McLean--Zinger]\label{thm:TMZ}
For every \(\omega\in\Symp^+(M,D)\), there is a path
\[
(\omega_t)_{0\leq t\leq1}\subset\Symp^+(M,D)
\]
with \(\omega_0=\omega\), \([\omega_t]=[\omega]\) for every \(t\), and such that
\(D\) admits an \(\omega_1\)-regularization.

Equivalently, by the Moser isotopy theorem, every SC symplectic divisor is
isotopic through SC symplectic divisors, inside a cohomologous family of
symplectic forms, to one admitting a regularization.
\end{theorem}

\begin{proof}
This is Theorem~7 and the conclusion immediately following it in
\cite{TMZ-CR}; the underlying normal-crossings regularization theorem is
Theorem~2.17 in \cite{TMZ-structures}.
\end{proof}

A regularization gives compatible tubular models for the strata \(D_I\). In
particular, for each \(I\) there is a direct sum of complex line bundles
\[
E_I=\bigoplus_{a\in I}L_a\longrightarrow D_I
\]
and a local identification of a disk subbundle of \(E_I\) with a neighborhood
of \(D_I\), under which the \(a\)-th branch is the zero locus of the
\(a\)-th fiber coordinate. These identifications satisfy the usual normal
crossings compatibility conditions on overlaps. A regularization also gives
complex line bundles \(\mathcal O_{\mathcal R}(D_a)\) with canonical sections
whose zero loci are \(D_a\); see \cite[Proposition~1.1 and \S2]{TMZ-structures}.

\section{The straightened almost K\"ahler normal form}

We need a slightly more explicit formulation of the usual existence theorem for
regularized almost K\"ahler data.

\begin{definition}\label{def:straight-chart}
An \emph{\(\mathcal R\)-straight chart} around \(p\in D_I\), where
\(I=\{1,\ldots,r\}\) after relabelling, consists of a neighborhood \(U\ni p\),
complex-valued functions \(z_1,\ldots,z_r\) on \(U\), a splitting
\[
TM|_U=H\oplus L_1\oplus\cdots\oplus L_r,
\]
and complex structures \(i_a\) on \(L_a\), such that
\begin{enumerate}
\item \(H=\ker dz_1\cap\cdots\cap\ker dz_r\);
\item \(z_a\) is a complex coordinate on the fibers of \(L_a\);
\item \(D_a\cap U=\{z_a=0\}\) for \(1\leq a\leq r\);
\item no branch \(D_b\) with \(b\notin I\) meets \(U\).
\end{enumerate}
An almost complex structure \(J\) on \(U\) is \emph{straight} in this chart if
it preserves \(H\) and every \(L_a\), and if \(J|_{L_a}=i_a\).
\end{definition}

\begin{proposition}[Straightened AK normal form]\label{prop:straight-AK}
Suppose that \(D\) admits an \(\omega\)-regularization \(\mathcal R\). Then
there are an \(\omega\)-tame almost complex structure \(J\) on \(M\) and an
\(\mathcal R\)-straight chart around every point of \(D\) in which \(J\) is
straight.
\end{proposition}

\begin{proof}
The argument is the connection-straight form of the construction of almost
K\"ahler data in \cite[\S2]{TMZ-structures} and \cite[\S2]{Tehrani-compact}.

Over a regularizing neighborhood of \(D_I\), identify \(M\) locally with an
open disk bundle in
\[
E_I=\bigoplus_{a\in I}L_a\longrightarrow D_I .
\]
Choose an \(\omega|_{D_I}\)-tame almost complex structure \(J_I\) on \(D_I\),
the given complex structures \(i_a\) on \(L_a\), and a horizontal distribution
\(H_I\). On the corresponding local trivialization, form the block-diagonal
almost complex structure
\[
J_I^{\mathrm{mod}}=J_I\oplus\bigoplus_{a\in I}i_a .
\]

The symplectic normal-form part of an \(\omega\)-regularization identifies
\(\omega\), after shrinking the disk bundle, with a form which is positive on
this horizontal--vertical splitting. More precisely, at the zero section the
horizontal part is \(\omega|_{D_I}\) and the vertical part is the positive
standard form on the symplectic normal lines; the remaining terms are
small on a sufficiently small disk bundle. Hence
\(J_I^{\mathrm{mod}}\) is \(\omega\)-tame after shrinking.

If \(z_a\) is the fiber coordinate associated with \(L_a\), then the horizontal
distribution is chosen to be \(\ker dz_a\) for every \(a\). Thus for
\(X\in H_I\),
\[
dz_a(J_I^{\mathrm{mod}}X)=0=i\,dz_a(X),
\]
while for \(v\in L_a\),
\[
dz_a(J_I^{\mathrm{mod}}v)=dz_a(i_av)=i\,dz_a(v).
\]
Consequently
\[
dz_a\circ J_I^{\mathrm{mod}}=i\,dz_a.
\]
Moreover,
\[
D_a\cap U=\{z_a=0\}.
\]

These local models are patched exactly as in the construction of
\(\AK(M,D)\). For a fixed \(\mathcal R\)-adapted chart, the conditions
\[
dz_a\circ J=i\,dz_a
\]
are affine in \(J\), and the \(\omega\)-tame condition is convex. The
compatibility conditions in the definition of a regularization identify the
indexed normal-line summands and their coordinate projections on overlaps.
Therefore convex combinations preserve the straight condition. A finite
partition of unity, followed by extension over \(M\), gives a global
\(\omega\)-tame \(J\). This is the same affine patching and contractibility
argument that shows that the projection from almost K\"ahler data to
regularization data is a weak homotopy equivalence; see
\cite[discussion after Theorem~7]{TMZ-CR} and
\cite[\S\S1--2]{Tehrani-compact}.
\end{proof}

\begin{remark}
It is important here that \(J\) be chosen in the connection-straight part of
the almost K\"ahler construction. Merely requiring \(D_I\) to be
\(J\)-holomorphic, or even requiring only the normal component of the
Nijenhuis tensor to vanish along \(D_I\), is the first-order condition commonly
used in relative pseudoholomorphic curve theory. Proposition
\ref{prop:straight-AK} records the additional local product normal form that
follows from regularized coordinates.
\end{remark}

\section{Proof of Theorem~\ref{thm:main}}

We now complete the proof.

\subsection*{Step 1: Deform to a regularized pair}

Apply Theorem~\ref{thm:TMZ}. This gives a cohomologous path
\((\omega_t)_{0\leq t\leq1}\) through symplectic forms for which \(D\) remains
an SC symplectic divisor, such that \(D\) admits an \(\omega_1\)-regularization.

Rename \(\omega_1\) as \(\omega\). Since the question permits the preliminary
cohomologous deformation, this causes no loss of generality.

\subsection*{Step 2: Choose straight almost K\"ahler data}

By Proposition~\ref{prop:straight-AK}, choose an \(\omega\)-tame almost complex
structure \(J\) and \(\mathcal R\)-straight charts in which \(J\) is straight.

Let \(p\in D\), and put
\[
I(p)=\{a:p\in D_a\},\qquad r=|I(p)|.
\]
After relabelling, \(I(p)=\{1,\ldots,r\}\). By the definition of an SC divisor,
there is a neighborhood of \(p\) meeting no branch \(D_b\) with \(b>r\).
Shrinking the straight chart if necessary, we may assume that it has this
property.

The straight chart supplies complex functions \(z_1,\ldots,z_r\) with
\[
dz_a\circ J=i\,dz_a,\qquad 1\leq a\leq r,
\]
and
\[
D\cap U
 =\bigcup_{a=1}^r(D_a\cap U)
 =\{z_1\cdots z_r=0\}.
\]

\subsection*{Step 3: Complete to a coordinate system}

Let
\[
(w_1,\ldots,w_{n-r})
\]
be any smooth complex coordinates on the stratum \(D_I\) near \(p\); no
holomorphicity is required of them. Pull them back to the regularizing tubular
neighborhood. The differentials
\[
dz_1,\ldots,dz_r,dw_1,\ldots,dw_{n-r}
\]
are complex-linearly independent at \(p\), because the first \(r\) differentials
are independent in the normal directions and the remaining ones restrict to a
coordinate coframe on \(D_I\). Therefore
\[
(z_1,\ldots,z_r,w_1,\ldots,w_{n-r})
\]
are smooth complex coordinates near \(p\). Relabel \(w_j\) as \(z_{r+j}\).
Then only \(z_1,\ldots,z_r\) are asserted to be \(J\)-holomorphic, and
\[
D\cap U=\{z_1\cdots z_r=0\}.
\]

If \(p\notin D\), choose a neighborhood \(U\subset M\setminus D\), take \(r=0\),
and choose arbitrary smooth complex coordinates on \(U\). This completes the
proof of Theorem~\ref{thm:main}. \qed

\begin{thebibliography}{9}

\bibitem{TMZ-CR}
M.~F. Tehrani, M.~McLean, and A.~Zinger,
\emph{Singularities and semistable degenerations for symplectic topology},
C.~R. Acad. Sci. Paris, Ser.~I \textbf{356} (2018), 420--432.
\href{https://www.numdam.org/item/10.1016/j.crma.2018.02.009.pdf}{numdam.org}.

% Required web-source link:
% [numdam.org](https://www.numdam.org/item/10.1016/j.crma.2018.02.009.pdf)

\bibitem{TMZ-structures}
M.~F. Tehrani, M.~McLean, and A.~Zinger,
\emph{Normal crossings singularities for symplectic topology: structures},
arXiv:2112.13125.
\href{https://ar5iv.labs.arxiv.org/html/2112.13125}{ar5iv.labs.arxiv.org}.

% Required web-source link:
% [ar5iv.labs.arxiv.org](https://ar5iv.labs.arxiv.org/html/2112.13125)

\bibitem{Tehrani-compact}
M.~F. Tehrani,
\emph{Pseudoholomorphic curves relative to a normal crossings symplectic
divisor: compactification},
arXiv:1710.00224.
\href{https://ar5iv.labs.arxiv.org/html/1710.00224}{ar5iv.labs.arxiv.org}.

% Required web-source link:
% [ar5iv.labs.arxiv.org](https://ar5iv.labs.arxiv.org/html/1710.00224)

\end{thebibliography}

\end{document}
```